\renewcommand\S{\mathbb{S}}

\begin{problem}[1]
  Let $\phi : V \to U$ be a linear map. Prove that $V/\Ker(\phi)$ is isomorphic to $\Im(\phi)$.
\end{problem}

\begin{solution}
  We want to show that $V/\Ker{\phi} \cong \Im(\phi)$. We will do this by showing that the map $\psi : V/\Ker{\phi} \to \Im(\phi)$ defined by $\psi(v + \Ker{\phi}) = \phi(v)$ is a well-defined linear isomorphism. The map $\psi(v + \Ker(\phi))$ is well-defined because if $v + \Ker(\phi) = w + \Ker(\phi)$, then $v - w \in \Ker(\phi)$, so $\phi(v - w) = 0$, which implies $\phi(v) = \phi(w)$. The map $\psi$ is linear because for any $v, w \in V$ and any scalar $c$, we have
  \begin{gather*}
    \psi((v + w) + \Ker(\phi)) = \psi(v + \Ker(\phi)) + \psi(w + \Ker(\phi)) = \phi(v) + \phi(w) = \phi(v + w) \\
    \psi(c(v + \Ker(\phi))) = c\psi(v + \Ker(\phi)) = c\phi(v) = \phi(cv)
  .\end{gather*}
  The map $\psi$ is injective because if $\psi(v + \Ker(\phi)) = 0$, then $\phi(v) = 0$, which implies $v \in \Ker(\phi)$, so $v + \Ker(\phi) = \Ker(\phi)$, which is the zero element in $V/\Ker(\phi)$. The map $\psi$ is surjective because for any $u \in \Im(\phi)$, there exists some $v \in V$ such that $\phi(v) = u$, so $\psi(v + \Ker(\phi)) = u$. Therefore, $\psi$ is a linear isomorphism, and we conclude that $V/\Ker(\phi) \cong \Im(\phi)$.
\end{solution}

\begin{problem}[2]
  Let $K$ be the simplicial complex shown in Figure 6.4 of the book, except with the direction of the arrow on the right flipped, so that $|K|$ is homeomorphic to the Klein bottle. Compute the Betti numbers of $K$. (You may use the fact that the Klein bottle is connected and non-orientable.)
\end{problem}

\begin{solution}
  Since $|K| \cong K'$, where $K'$ is the Klein bottle, we can compute the Betti numbers of $K$ by computing the Betti numbers of $K'$. For $\beta_0$, since $K'$ is connected, it has one connected component. Therefore, $\beta_0 = 1$. For $\beta_1$, since $K'$ is non-orientable, it has one independent cycle. Therefore, $\beta_1 = 1$. For $\beta_2$, since $K'$ is non-orientable, it has no 2-dimensional holes. Thus, $\beta_2 = 0$. Thus, we have $\beta_0 = 1$, $\beta_1 = 1$, and $\beta_2 = 0$.
\end{solution}

\begin{problem}[3]
  Let $\Gamma$ be a tree, that is, a connected 1-dimensional simplicial complex with n vertices and $n - 1$ edges for some $n > 0$. Compute the Betti numbers of $\Gamma$.
\end{problem}

\begin{solution}
  For $\beta_1$, since $\Gamma$ is a tree, it has no cycles. Therefore, $\beta_1 = 0$. For $\beta_0$, since $\Gamma$ is connected, it has one connected component. Therefore, $\beta_0 = 1$. For $\beta_2$, since $\Gamma$ is a 1-dimensional simplicial complex, it has no 2-dimensional holes. Thus, $\beta_2 = 0$. Thus, we have $\beta_0 = 1$, $\beta_1 = 0$, and $\beta_2 = 0$.
\end{solution}

\begin{problem}[4]
  Let $\Gamma$ be the complete graph on 4 vertices, meaning that it is a connected 1-dimensional simplicial complex with 4 vertices and 6 edges. Compute the Betti numbers of $\Gamma$.
\end{problem}

\begin{solution}
  For $\beta_0$, since $\Gamma$ is a complete graph, it is connected. Therefore, $\beta_0 = 1$. For $\beta_1$, the complete graph on 4 vertices and 6 edges. So, it has $6 - 3 + 1 = 3$. Therefore, $\beta_1 = 3$. For $\beta_2$, since $\Gamma$ is a 1-dimensional simplicial complex, it has no 2-dimensional holes. Thus, $\beta_2 = 0$. Thus, we have $\beta_0 = 1$, $\beta_1 = 3$, and $\beta_2 = 0$.
\end{solution}

\begin{problem}[5]
  Let $\Delta_1$ be the simplex in $\R^3$ with vertices $\{(0, 0, 0), (1, 0, 0), (0, 1, 0), (0, 0, 1)\}$. Let $\Delta_2$ be the simplex in $\R^3$ with vertices $\{(0, 0, 0), (1, 0, 0), (0, 1, 0), (0, 0, -1)\}$. Let $K$ be the simplicial complex consisting of all proper faces of $\Delta_1$ along with all proper faces of $\Delta_2$. Thus the realization $|K|$ is homeomorphic to two copies of $\S^2$, glued together along a copy of $D^2$. Compute the Betti numbers of $K$.
\end{problem}

\begin{solution}
  For $\beta_0$, since $|K|$ is homeomorphic to two copies of $\S^2$ glued together along a copy of $D^2$, it is connected. Therefore, $\beta_0 = 1$. For $\beta_1$, since $|K|$ is homeomorphic to two copies of $\S^2$ glued together along a copy of $D^2$, it has no independent cycles. Therefore, $\beta_1 = 0$. For $\beta_2$, since $|K|$ is homeomorphic to two copies of $\S^2$ glued together along a copy of $D^2$, it has one independent 2-dimensional hole, so $\beta_2 = 2$. Thus, we have $\beta_0 = 1$, $\beta_1 = 0$, and $\beta_2 = 2$.
\end{solution}
